% !TEX root = ../monograph.tex
\chapter{Large fields and local Gaussian integration}
\label{ch06:large-fields}

A step integrates every configuration. The usefulness of
small-field coordinates therefore depends on a bound for their complement.
This chapter separates two mechanisms. Large plaquette fields and
large deviations of coarse links from their averages, both indexed by a large-field index triple (a set of fine plaquettes, a set of coarse edges and a set of coarse sites, defined below), retain a deterministic positive energy.
Large fluctuation-field events, the failures of smooth local cutoffs in the remaining coordinates, are controlled
by simultaneous Gaussian tail estimates. These mechanisms use different
estimates and must be combined under the same integral. Ba{\l}aban makes the same split, into large plaquette fields and large fluctuation fields~\cite{Balaban1988Expansions}; he adds the location of a large fluctuation field to the large-field region, whereas here the large-field set is not enlarged and these events are bounded by a joint Gaussian tail estimate.

We work at two consecutive scales, the fine scale $m+1$ and the coarse scale $m$. The algebraic decompositions below
are exact. The energy and differentiated integral estimates use
uniform bounds for the chosen path components and for the non-flat matrices of a sector (a term of the decomposition below with fixed discrete data).
We construct these bounds for the same local integral below; the flat
Gaussian identities alone would not supply them.

\section{Large-field index triples and the local energy}

Fix a stage, a block ratio $L$, the penalty parameter $b>0$, and the fine and coarse
path data of \eqref{ch04:path-action}. Throughout this chapter distances
are coarse lattice distances, except that path-length truncations use true fine
length. Write
\[
 t_f=t_{m+1},\qquad t_c=t_m,\qquad B=B_m,\qquad
 R_*=\lfloor B/100000\rfloor,\qquad
 t_{\rm hi}=\widehat\beta_{m+1}^{-49/100}.
\]
The thresholds $t_m=\widehat\beta_m^{-1/2}M_m^{\rm cut}$ and
$B_m=\lceil(\log\widehat\beta_m)^2\rceil$ were defined in
\eqref{ch04:thresholds}. At fixed $L$, as $\widehat\beta_m\to\infty$,
\begin{equation}
 \frac{t_f}{t_c}\longrightarrow1,\qquad
 B^n\frac{t_f}{t_{\rm hi}}\longrightarrow0,\qquad
 B^n\frac{t_{\rm hi}^2}{t_f}\longrightarrow0
 \quad(n\text{ fixed}).
 \label{ch06:scales}
\end{equation}
The coarse torus period is at least $100(B+1)$. All constants below are
independent of that period.

For a configuration pair $(U,V)$ and the value $Q(U)$ of the averaging operation of
Chapter~\ref{ch04:blocking}, put
\[
 r_p^f=\chi(U_p),\qquad r_p^c=\chi(V_p),\qquad
 n_e=\chi(V_eQ_e(U)^{-1}).
\]
Here $\chi$ is the normalized Hilbert--Schmidt distance from the identity, and $n_e$ is called the \emph{deviation} of the coarse link $V_e$ from its average $Q_e(U)$. A \emph{large-field index triple}, for short an \emph{index triple},
$A=(A_f,A_n,A_c)$ consists of fine plaquettes, coarse edges, and
coarse sites. Its size is the sum of the three cardinalities. Its location
contains the projected corners of its plaquettes, the endpoints of its
edges, and its sites. Write $A^{[r]}$ for the coarse $r$-neighborhood of
that location. Fix $W>0$. An index triple is consistent if, at all plaquettes and edges meeting
$A^{[5B]}$,
\begin{align*}
 &r_p^f\ge t_f/4\quad(p\in A_f),&
 &r_p^f\le t_f/2\quad(p\notin A_f),\\
 &n_e\ge Wt_f\quad(e\in A_n),&
 &n_e\le2Wt_f\quad(e\notin A_n),
\end{align*}
and $A_c$ is exactly the set of coarse sites incident to plaquettes with
$r_p^c>t_c$ in that neighborhood. Every chosen value of the auxiliary index of Chapter~\ref{ch04:blocking} for a local gauge transformation must have positive
probability weight. The overlaps in the first two threshold pairs permit a partition
with smooth cutoffs; consistency does not prescribe a unique index triple.

For the exact integral use a global partition of unity before forming components.
Choose a smooth function $\vartheta_f$ equal to zero below $t_f/4$
and to one above $t_f/2$, with values in $[0,1]$, and an analogous
$\vartheta_n$ for the thresholds $Wt_f,2Wt_f$. Give a set $A_f$ of fine plaquettes
the weight
\[
 \prod_{p\in A_f}\vartheta_f(r_p^f)
 \prod_{p\notin A_f}(1-\vartheta_f(r_p^f)),
\]
and use the analogous product for the sets $A_n$ of coarse edges. Summing over all
pairs $(A_f,A_n)$ gives one at every configuration. Take $A_c$ to be the
coarse large-field set of the configuration. Each index triple of nonzero weight obeys
the consistency inequalities at every plaquette and edge, not merely within $5B$.
The weaker local definition above suffices for the bound on the local energy $\mathcal E_A$ defined below;
later component factors retain the restrictions of the global partition of unity
on their larger dependence sets. This refinement adds local
eligibility conditions without discarding any part of the original integral.

A \emph{sector} is a term of this decomposition with fixed discrete data: it fixes the index triple, its assignment to connected
components of its neighborhoods, a rooted spanning tree in every block,
and the value of the auxiliary index. Tree links are set to the identity.
These choices fix the coordinates and supports before differentiation.
Let
\[
 D=\{p\in A_f:r_p^f>t_{\rm hi}\}\cup\{e\in A_n:n_e>t_{\rm hi}\}.
\]
Let $H$ be the coarse plaquettes based at distance at most $4B/5$ from $D$. Local
counting gives
\begin{equation}
 |H|\le K B^4|D|,\qquad |A^{[2B]}|\le K B^4|A|.
 \label{ch06:counts}
\end{equation}
The constants may depend on the fixed stage and ratio.

\subsection{An exact completion of squares}

Let $F(g)=P_{\gLie}(g-g^{-1})/2$, with the orthogonal projection taken in
the fixed trace metric. The plaquette and path components of the curvature are grouped as
\[
 X_f(p)=\bigl(\sqrt{c_{\rm dir,f}}F(U_p),L_{f,R_*}(U,p)\bigr),
 \quad
 X_c(p)=\bigl(\sqrt{c_{\rm dir,c}}F(V_p),L_{c,R_*}(V,p)\bigr).
\]
The $L_{*,R_*}$ are the path sums of \eqref{ch04:path-action}, truncated
at true fine length $R_*$. Let $\chi_p^{\rm rel}$ denote the indicator of the relaxed
regularity condition \eqref{ch04:relaxed-test}. Define coarse components, whose plaquette part carries the factor $\alpha_p$,
and local action densities by
\begin{align*}
 v_c(p)&=\bigl(\alpha_p\sqrt{c_{\rm dir,c}}F(V_p),
                   \chi_p^{\rm rel}(V)L_{c,R_*}(V,p)\bigr),
 &\alpha_p&=\min(1,t_c/r_p^c),\\
 K_f(p)&=\operatorname{clip}_{t_f}(r_p^f)
                  +\tfrac12\chi_p^{\rm rel}(U)|L_{f,R_*}(U,p)|^2,\\
 K_c(p)&=\operatorname{clip}_{t_c}(r_p^c)
                  +\tfrac12\chi_p^{\rm rel}(V)|L_{c,R_*}(V,p)|^2.
\end{align*}
Set $\alpha_p=1$ when $r_p^c=0$. Here $\operatorname{clip}_t$ is the Borel quadratic/linear
function in \eqref{ch04:clip}, with its respective scale coefficients.
In particular $|v_c(p)|\le Ct_f$.

For clarity the square-completion input can be stated without choosing a
particular factorization of the metrics. Suppose flat maps for the fine curvature, the penalty term, and the
complement, $M_f,M_n,N$, satisfy
\begin{equation}
 M_f^*M_f+M_n^*M_n+N^*N=I,
 \qquad M_f^*X_f^{\rm lin}+M_n^*Z^{\rm lin}=X_c^{\rm lin}.
 \label{ch06:flat-identities}
\end{equation}
The superscript $\mathrm{lin}$ denotes linearization at the flat field.
Here $a$ and $v$ are the fine and coarse link variations, respectively,
and $Z^{\rm lin}=\sqrt b(v-\mathcal Q a)$ is their discrepancy under the
linearized averaging operator $\mathcal Q$ of \eqref{ch05:actual-Q}. For the fixed quadratic data these identities follow from the
identities \eqref{ch05:physical-transports} and \eqref{ch05:positive-defect} of Chapter~\ref{ch05:gaussian}.
Here is their realization through a factorization of the metrics, including the exact factorization hypothesis.
Let $C_\nu$ map curvature to the plaquette and path components used at
scale $\nu\in\{f,c\}$, and require
$C_\nu^*C_\nu=\mathsf H_\nu$, where $\mathsf H_\nu$ is the chosen
modified curvature metric. This is the quadratic coefficient condition
of \eqref{ch05:actual-hessian}; it is not a conclusion for an arbitrary
path action. Put $\mathsf B_\nu=\mathsf H_\nu^{-1}$ and
$P_c=C_c\mathsf B_cC_c^*$. The positive plaquette component makes $C_c$
injective. Thus $P_c=C_c(C_c^*C_c)^{-1}C_c^*$ is the orthogonal
projection onto the range of $C_c$. Equation
\eqref{ch05:positive-defect} supplies an exponentially local operator $K$
with
\[
 K^*K=\mathsf B_c-T^*\mathsf B_fT-b^{-1}d_{1,c}d_{1,c}^*.
\]
Choose
\[
 M_f=C_f\mathsf B_fTC_c^*,\qquad
 M_n=b^{-1/2}d_{1,c}^*C_c^*,\qquad
 N=\begin{pmatrix}I-P_c\\KC_c^*\end{pmatrix}.
\]
Direct multiplication gives the first identity in
\eqref{ch06:flat-identities}. For the second, use
$X_f^{\rm lin}=C_fd_{1,f}a$,
$X_c^{\rm lin}=C_cd_{1,c}v$, and
$T^*d_{1,f}=d_{1,c}\mathcal Q$. The last identity is
\eqref{ch05:physical-transports} with the linearization $\mathcal Q$ of the averaging operation in
\eqref{ch05:actual-Q}: its extra coarse gradient has zero curl.
The metric modification does not change these identities on curls.
Thus both identities \eqref{ch06:flat-identities} hold for the same action components satisfying the factorization hypothesis and
the same linearization $\mathcal Q$. The path construction used here imposes this factorization:
in Section~\ref{ch07:classical} it chooses
$C_c=(\sqrt{c_{\rm dir}}I,\rho)^T$ with
$\rho^*\rho=\mathsf H_c-c_{\rm dir}I$, and realizes these very components
by parallel-transported plaquettes. Cubic and quartic matching terms do not
change their Hessian, while the $\beta$-independent quadratic term of the path action is not part
of the classical metric. The step in Theorem~\ref{ch09:step-theorem}
uses precisely this construction. At the Wilson initialization (the Wilson density at the ultraviolet cutoff) the plaquette
and additional components have coefficients $\sqrt{c_{\rm dir}}$ and
$\sqrt{1-c_{\rm dir}}$, so their Gram matrix is $I$; the
fine metric at the ultraviolet cutoff is unmodified ($\zeta_f=0$). The first coarse modification
and every later one obey \eqref{ch05:positive-defect}. Thus the factorization
is an invariant of this construction, although it remains a
necessary premise when applying the estimate to other path actions.
Parallel-transport the coefficients
along their fixed coordinate contours using the pull-back $\operatorname{rep}V$ of the coarse field,
and discard contours of fine length greater than $R_*$. Denote the resulting
maps by $\Phi,\Phi_n,\mathcal N$ and put
$Z(e)=\sqrt b F(V_eQ_e^{-1})$. Their adjoints below are the adjoints of the
parallel-transported maps.

The vectors $\Phi v_c$ and $\Phi_n v_c$ are the parallel-transported linear
comparison terms. Throughout this chapter, a comparison term means one
of these specified vectors, not an independently chosen approximation.

Define six local terms:
\begin{align}
 E_f&=K_f+\tfrac12|\Phi v_c|^2-\langle\Phi v_c,X_f\rangle,\notag\\
 E_n&=\tfrac b2 n_e^2+\tfrac12|\Phi_n v_c|^2
                                      -\langle\Phi_n v_c,Z\rangle,\notag\\
 E_c&=\langle v_c,X_c\rangle-\tfrac12|v_c|^2-K_c,\notag\\
 E_N&=\tfrac12|\mathcal N v_c|^2,\notag\\
 E_L&=\langle v_c,\Phi^*X_f+\Phi_n^*Z-X_c\rangle,\notag\\
 E_D&=\tfrac12\langle v_c,
       (I-\Phi^*\Phi-\Phi_n^*\Phi_n-\mathcal N^*\mathcal N)v_c\rangle.
 \label{ch06:six-terms}
\end{align}
Every expression is evaluated at its corresponding plaquette, edge, or
complement-component base point. On summing over the whole torus, the cross terms
and squares cancel exactly:
\begin{equation}
 \sum(E_f+E_n+E_c+E_N+E_L+E_D)
       =\sum K_f+\frac b2\sum_e n_e^2-\sum K_c.
 \label{ch06:global-energy}
\end{equation}
This equality is algebraic and does not assume flatness. The local energy
$\mathcal E_A$ is instead the sum of the six terms at base points in
$A^{[2B]}$. Cancellation across its boundary is unavailable. The two terms $E_L,E_D$, which are not of definite sign,
must therefore be estimated locally.

\subsection{Direct bounds and the signed comparison}

Assume the parallel-transported maps have absolute row and column sums bounded by
$CL^{-2}$ and $CL^2$ for $\Phi$, respectively, and by $C$ for
$\Phi_n,\mathcal N$. Bounds involving $L^{-2}$ below require $C$ fixed
before choosing $L$. Put
\[
 R_f=\sum_{p\in A_f:r_p^f>t_f}r_p^f,
 \qquad R_n=\sum_{e\in A_n}n_e.
\]
The fine plaquette term, when $r_p^f\le t_f$, is at least
$\tfrac12(\sqrt{c_{\rm dir,f}}r_p^f-|\Phi v_c(p,0)|)^2$.
Since $|\Phi v_c|\le CL^{-2}t_f$, a plaquette of $A_f$ in this range
contributes at least $c_{\rm dir,f}t_f^2/128$ for large $L$.
Above $t_f$ the linear clip retains a multiple of $t_fr_p^f$.

There is one additional loss: a path component whose relaxed regularity condition fails is no
longer part of a square. Such a component based at $p$ has a witness $q$ with
\[
 r_q^f>t_f\exp(\operatorname{dist}(p,q)/R^{\rm rel}).
\]
For a fixed witness the number of possible bases is at most
$C(1+R^{\rm rel}\log(r_q^f/t_f))^4\le C r_q^f/t_f$.
All witnesses lie within $B/50$ for sufficiently large reference coupling.
They belong to $A_f$ by consistency. Split path endpoints
into those with distance from the identity at most $t_f$ and the rest. The row bound handles
the former, and the column bound handles the latter, giving
\[
 \sum_{\text{omitted components }p}|L_{f,R_*}(U,p)|\le C R_f.
\]
Multiplication by $|\Phi v_c|\le CL^{-2}t_f$ absorbs this loss into the
linear clip after increasing $L$. Consequently
\begin{equation}
 \sum E_f\ge\frac{c_{\rm lin,f}}2 t_fR_f
  +\frac{c_{\rm dir,f}}{128}t_f^2
                         |\{p\in A_f:r_p^f\le t_f\}|.
 \label{ch06:fine-bound}
\end{equation}
The square of the deviation similarly gives, for sufficiently large $W$ after $L$,
\begin{equation}
 \sum E_n\ge\frac b8\sum_{e\in A_n}n_e^2.
 \label{ch06:mismatch-bound}
\end{equation}

The coarse path contribution cancels exactly in $E_c$. The plaquette part is
$c_{\rm dir,c}(\alpha_p-\alpha_p^2/2)|F(V_p)|^2
-\operatorname{clip}_{t_c}(r_p^c)$.
The squared normal displacement from the tangent space is $O(\chi(g)^4)$,
so $|F(g)|^2\ge\chi(g)^2-C\chi(g)^4$ near the identity. Choose fixed
$\delta>0$ so that also $|F(g)|^2\ge3\chi(g)^2/4$ below $\delta$.
The clip coefficients satisfy $c_{\rm dir,c}>16c_{\rm lin,c}$, and hence
\begin{equation}
 E_c(p)\ge
 \begin{cases}
 -Ct_c^4,&r_p^c\le t_c,\\
 c_{\rm dir,c}t_cr_p^c/4,&t_c<r_p^c\le\delta,\\
 -c_{\rm lin,c}t_cr_p^c,&r_p^c>\delta.
 \end{cases}
 \label{ch06:coarse-bound}
\end{equation}

The remaining ingredient is genuinely about the chosen nonabelian paths
and averaging operation. Its quantitative form is recorded explicitly.
\begin{proposition}[Local signed comparison from compatible components]
\label{ch06:signed-input}
Assume \eqref{ch06:flat-identities}, the stated row and column
bounds for the parallel-transported maps, and a common positive exponential bound for the uncut path
kernels, with their coefficients satisfying the factorization hypothesis for the curvature components used
in those identities. Then, on every consistent sector for the averaging operation
of Chapter~\ref{ch04:blocking}, there are $C,K$ fixed before $L,W$ and
$Z=Z_{j,L,W}$ such that
\begin{align}
 \sum_H(E_L+E_D)&\ge
 c_{\rm lin,c}t_c\sum_{r_p^c>\delta}r_p^c
       -Ct_f(L^{-2}R_f+R_n)-Zt_f^2B^4|D|,
 \label{ch06:affected}\\
 \sum_{A^{[2B]}\setminus H}(E_L+E_D)&\ge
                                  -\frac a4t_f^2|A|.
 \label{ch06:unaffected}
\end{align}
Every coarse plaquette with $r_p^c>\delta$ lies in $H$ and their number
is at most $K|D|$. Here
\[
 a=\min\{c_{\rm lin,f}/8,c_{\rm dir,f}/128,b/32,c_{\rm dir,c}/64\}.
\]
\end{proposition}

\begin{proof}
In the lifted macro-box
of radius $(j+7)L$, let $\mathcal P_y(U)=\sum_qm_y(q)r_q^f$ be the sum of
fine plaquette distances from the identity with the contour multiplicities $m_y(q)$. The required comparison
is
\begin{equation}
 r_p^c\le\sum_{e\in\partial p}n_e+C_jL^{-2}\mathcal P_y(U),
 \qquad \sup_q\sum_y m_y(q)\le K_j.
 \label{ch06:macro-comparison}
\end{equation}
This estimate does hold for the precise branches of
Chapter~\ref{ch04:blocking}, as the following argument shows. It is a
bound linear in the plaquette sum $\mathcal P_y(U)$; comparison of squared energies of the local gauge transformations alone would give
an inadequate square-root bound.

The estimate for the averaging operation, linear in the plaquette sum, is Lemma~\ref{ch04:predictor-comparison}; its comparison uses the same axial gauge on the macro-box and the same rooted contours as here.

The same bound localizes coarse plaquettes with $r_p^c>\delta$.
If every variable $r_q^f$ or $n_e$ in their fixed macro-box is at most
$t_{\rm hi}$, \eqref{ch06:macro-comparison} bounds $r_p^c$ by
$C_{j,L}t_{\rm hi}<\delta$ after increasing the lower bound on the reference coupling.
Thus each such plaquette has a plaquette or edge of $D$ in that macro-box. The fixed
macro-box fits inside the $B$-neighborhood used to define $H$, and bounded
multiplicity of the macro-boxes gives the asserted count $K|D|$.
Separating the small-distance contribution from the large-field contribution $\Delta_p$ gives
$r_p^c\le C_{j,L,W}t_f+\Delta_p$ with
$\sum_p\Delta_p\le C(L^{-2}R_f+R_n)$.
The constant $C$ in the last bound must remain independent of $L,W$.
Applying the row and column bounds for the parallel-transported maps then gives
\eqref{ch06:affected}; the positive displayed term is obtained by absorbing
its upper bound into the enlarged loss constants. It is not a new
positivity assertion for $E_L$.

Outside $H$, all fine plaquette distances from the identity and deviations $n_e$ in the
surrounding fixed macro-boxes are at most $t_{\rm hi}$. The variables at
plaquettes and edges outside $A_f$, $A_n$ obey this too, since $2Wt_f\le t_{\rm hi}$ eventually. The
comparison \eqref{ch06:macro-comparison} bounds coarse plaquette distances from the identity there by
$C_{j,L,W}t_{\rm hi}$. Choose a coarse axial gauge on a box of radius
$B/2$. Its sum of bond distances from the identity is bounded by $C_{j,L,W}B^5t_{\rm hi}$.
To compare a fine link with the corresponding link of $\operatorname{rep}V$, join its two
endpoints to their block roots along the tree. All contour variables along the tree
are the identity. The resulting contour lies in two neighboring blocks.
In an axial gauge on the macro-box its distance from the identity is bounded by the length of the contour times
the sum of bond distances from the identity in the macro-box. Comparing a crossing contour with the value of the averaging operation
adds the errors of the local gauge transformations at the two endpoints, already bounded in
Lemma~\ref{ch04:predictor-comparison}, and then the deviation $n_e$.
Consequently, with a constant depending on the fixed tree family and $L$,
\[
 \chi\bigl(U_e(\operatorname{rep}V)_e^{-1}\bigr)
       \le C_{j,L}\mathcal P_y(U)+n_{e_c}
\]
for crossing edges; omit the deviation for internal edges. Any tree in a
fixed block has uniformly bounded path length (at most $L^4-1$), which
suffices here because this constant may depend on $L$. This is not a
bound on $\Delta_p$ with a constant independent of $L$. The parent lift of the coarse gauge preserves
these distances between corresponding links. It therefore places both fine and coarse
logarithms in a common chart, throughout radius $2B/5$, with bound
$\delta_{\rm loc}=C_{j,L,W}B^{10}t_{\rm hi}$.

The derivative of the rooted averaging operation also matches the flat
identity. Near the identity all branches of the local gauge transformations and of the averaging operation chosen by the auxiliary index
are minimizing; the redundant values of the auxiliary index do not change their values. The
finite-patch minimizing local gauge transformation is smooth there by the same nondegenerate
implicit equation used in Lemma~\ref{ch04:frame-lemma}. Write
$h_{y,z}(\exp(\epsilon a))=1+\epsilon H_{y,z}(a)+O(\epsilon^2)$ and
$\overline H_y=\sum_z w_z^{\rm av}H_{y,z}$. Differentiating
the transporter \eqref{ch04:connector}, and then the averaging operation applied to it, gives
\[
 (\mathcal Qa)_{y,\mu}
    =(\mathcal Aa)_{y,\mu}+\overline H_y(a)
                                -\overline H_{y+e_\mu}(a),
\]
where $\mathcal A$ averages the straight length-$L$ link sums with the
same weights. Thus $\mathcal Q-\mathcal A$ is a coarse gradient. Its curl
vanishes exactly. This verifies the identification of the linearized averaging operation needed in
\eqref{ch06:flat-identities}. Their realization through the factorization of the metrics was derived
above from \eqref{ch05:physical-transports} and
\eqref{ch05:positive-defect}; the incoming path action must still satisfy
the exact quadratic coefficient condition in
\eqref{ch05:actual-hessian}. Second-order
Taylor bounds are uniform in the period because each value of the averaging operation reads a
fixed finite patch at fixed $j,L$.

Set
$\varepsilon_{\rm loc}=(1+R_*)\delta_{\rm loc}+e^{-\tau R_*}$,
where $\tau>0$ is a common path decay rate. Subtracting the flat identities
\eqref{ch06:flat-identities} before estimating terms gives the required
bounds
\begin{align*}
 |\Phi^*X_f+\Phi_n^*Z-X_c|&\le C\delta_{\rm loc}\varepsilon_{\rm loc},\\
 |(I-\Phi^*\Phi-\Phi_n^*\Phi_n-\mathcal N^*\mathcal N)v_c|
                                    &\le Ct_f\varepsilon_{\rm loc}.
\end{align*}
Indeed the contour variable of a contour of length $s$ differs from its flat value by at most
$Cs\delta_{\rm loc}$; the omitted tails have size $Ce^{-\tau R_*}$.
The count \eqref{ch06:counts} and the limits \eqref{ch06:scales} yield
\eqref{ch06:unaffected}. This completes the verification of Proposition~\ref{ch06:signed-input}.
It uses the flat identities, with their factorization, for the chosen path family, not merely
a decay estimate for a flat Fourier kernel.
\end{proof}

\begin{proposition}[Energy bound under the local comparisons]
\label{ch06:energy}
Under the hypotheses of Proposition~\ref{ch06:signed-input}, choose
$L$ sufficiently large, then $W$ sufficiently large, then the reference
coupling sufficiently large. Uniformly in the period and consistent sector,
\[
 \mathcal E_A(U,V)\ge\kappa t_f^2|A|,\qquad \kappa=a/2>0.
\]
The lower bound on the reference coupling may depend on $j,L,W$; $a$ is fixed by the coefficients $c_{\rm dir}$, $c_{\rm lin}$ of the
clip functions and by $b$.
\end{proposition}
\begin{proof}
The bounds \eqref{ch06:fine-bound} and \eqref{ch06:mismatch-bound} leave
$a t_f^2$ per plaquette of $A_f$ or edge of $A_n$, $a t_ft_{\rm hi}$ per element
of $D$, and enough linear surplus to absorb
$Ct_f(L^{-2}R_f+R_n)$. The choices are in that order: increase $L$ for
the term $L^{-2}R_f$, and $W$ for the term $R_n$.

Assign each coarse large-field site to one incident plaquette with distance from the identity above
$t_c$. A plaquette has at most four such sites. Since $t_f\le2t_c$
eventually, the middle range of \eqref{ch06:coarse-bound} contributes
$4at_f^2$ per plaquette and therefore $at_f^2$ per assigned site. The
exceptional plaquettes with distance from the identity above $\delta$ number at most $K|D|$;
their negative energies cancel the positive term in
\eqref{ch06:affected}. Their missing site contributions amount to at most
$4aKt_f^2|D|$. Also $E_N\ge0$. The total remaining losses are bounded by
\[
 (4aK+ZB^4)t_f^2|D|+Ct_c^4KB^4|A|+\frac a4t_f^2|A|.
\]
Increase the lower bound on the reference coupling until
$(4aK+ZB^4)t_f^2\le at_ft_{\rm hi}$ and
$Ct_c^4KB^4\le at_f^2/4$. Equation~\eqref{ch06:scales} justifies both
choices. The first loss is absorbed by the positive energy contribution retained for $D$, leaving
at least $at_f^2|A|/2$.
\end{proof}

\section{The Gaussian of a sector at a non-flat background}
\label{ch06:sector}

Fix the sector and the conditioning variables $Y$: the non-tree links that are held fixed in the Gaussian integral and integrated afterwards, called \emph{frozen} below. (The word conditioning is Ba{\l}aban's~\cite{Balaban1988Cluster}; in~\cite{Balaban1988Expansions} the integrations over such variables are left to later steps, whereas here they are carried out in the same step.) Tree links are not integration
variables. \emph{Active} coordinates, the integration variables, are logarithms (Lie-algebra coordinates) of the remaining non-tree links
outside $A^{[10B]}$. Linearize the retained square
terms at the pull-back $\operatorname{rep}V$ of the coarse field, using the chosen averaging operation
and keeping the cutoffs fixed. Collect the resulting affine expressions
in a single vector $\mathsf D u+b_Y$, where $u$ contains the active
coordinates and $b_Y$ is the value at $u=0$. Their quadratic energy is
\begin{equation}
 Q_Y(u)=\tfrac12|\mathsf D u+b_Y|^2,\qquad
 \mathsf A=\mathsf D^*\mathsf D,\qquad f_Y=\mathsf D^*b_Y.
 \label{ch06:actual-matrix}
\end{equation}
If hypercubic orientations are averaged, this affine map uses the square roots
of their weights. A component is retained when its \emph{link support}, the fixed set of links from which its expression is built, meets an active
edge. This rule differs from retaining components whose base points are a specified
distance from the location of $A$. A comparison between the corresponding
matrices requires a proof. Identity coordinates may be adjoined for
inactive variables, with their trivial integrals normalized to one.

The link support contains the links on which the averaging operation depends and every endpoint of a retained
path; here $r_{\rm pred}$ is the largest coarse distance at which the averaging operation reads a link. Its coarse radius is at most $R_*+5+r_{\rm pred}\le B/50$ for large
coupling. Thus the link support of a retained component lies in the small-field region when its
active edge is outside $A^{[10B]}$. This geometric
observation permits local estimates at the background field $\operatorname{rep}V$; it does not by
itself prove a lower bound on $\mathsf A$.

\subsection{Transferring a local estimate in the block axial gauge to the non-flat components}

The gap needed below can be obtained without identifying non-flat and flat
matrices. The geometric comparison just proved implies a local inequality in the
specified block axial gauge:
\begin{equation}
 |a_e|\le C_L\sum_{p\in\mathcal P_e}|(d_fa)_p|
                    +|\mathcal Q a(e_c)|.
 \label{ch06:tree-control}
\end{equation}
Here $\mathcal P_e$ is the set of fine plaquettes of a fixed macro-box,
$e_c$ is the coarse edge in the direction of $e$, and $\mathcal Q$ is
the derivative of the averaging operation at the identity. The inequality
is for arrays vanishing on the chosen tree, extended by zero on frozen
edges. For internal edges the last term may be omitted. Its nonlinear
origin is worth specifying: join the endpoints of $e$ to their block
roots along the same tree, compare the resulting contour variable with the value of the
averaging operation in a common axial gauge on the macro-box, and bound the difference by the
plaquette sum of the macro-box. Differentiating this comparison at the identity
gives \eqref{ch06:tree-control}. The contour constants and endpoint bounds are those proved above for the
same trees and the same values of the auxiliary index. No coordinate change to an unrelated tree is made.

The Poincar\'e-type inequality \eqref{ch05:tree-comparison} applies to these
coordinates: the pull-back $\operatorname{rep}V$ is varied by left multiplication on non-tree bonds,
tree coordinates are zero, and the averaging operation is the rooted one, with the linearization $\mathcal Q$
of \eqref{ch05:actual-Q}. Thus the retained plaquette components and the retained components containing the averaging operation, for short \emph{averaging components},
are exactly its derivatives, multiplied by their fixed positive
square weights. Each such derivative differs, in one common gauge on its
macro-box, from its flat derivative by at most $\varepsilon$ times the
sum of input norms on its fixed link support. This follows from continuity of
the finite-patch derivative on the minimizing branch once clearance and
small coarse plaquettes put that entire patch in a small common gauge.
Insert the comparison in \eqref{ch06:tree-control}, square, and sum over
active edges. Fixed multiplicities of the link supports give
\[
 \|a\|_2^2\le C_L\sum_p|D_p a|^2+C_L\sum_e|D_e^{\rm pred}a|^2
                         +C_L\varepsilon^2\|a\|_2^2.
\]
All norms here are unnormalized $\ell^2$ norms, without division by volume. The choice of retained
components introduces no error: an omitted component has a link support disjoint
from the active coordinates, hence annihilates their zero extension.
For $C_L\varepsilon^2\le1/2$ and positive lower plaquette and penalty
coefficients, this proves
$\mathsf D^*\mathsf D\ge c_LI$. The other components add nonnegative squares.
The gauge may depend on the edge at which the local inequality is applied.
It is fixed during differentiation, and invariance of the norm on $\gLie$
permits the subsequent sum. The resulting lower bound $c_L$ may depend
on $L$. The constant in the block comparison
\eqref{ch06:macro-comparison} must already be fixed when $L$ is chosen.

\subsection{A localized approximate inverse}

For an edge coordinate $i$, let $E_i$ be its two projected endpoints.
Two useful symmetric windows are
\begin{align*}
 \Omega_i^{\rm sec}&=\{i':\max_{x\in E_i,y\in E_{i'}}
                                      \operatorname{dist}(x,y)\le B/20\},\\
 \Omega_i^{\rm loc}&=\{i':\min_{x\in E_i,y\in E_{i'}}
                                      \operatorname{dist}(x,y)\le B\}.
\end{align*}
Restrict them to the active coordinate set. Both contain a ball of radius
$B/60$ about $i$ for large $B$. Use one prescription throughout a given
integral and define
\begin{equation}
 P_{i'i}=\mathbf1_{i'\in\Omega_i}
                    (\mathsf A[\Omega_i,\Omega_i]^{-1})_{i'i},
 \qquad \Gamma=\tfrac12(P+P^*),\qquad m_Y=-P^*f_Y.
 \label{ch06:windows-def}
\end{equation}
Here $\mathsf A[\Omega_i,\Omega_i]$ denotes the principal submatrix on
$\Omega_i$: to obtain column $i$ of $P$, solve the restricted system with
right-hand side the $i$th coordinate vector and extend the solution by
zero. The comparison below shows that the finite-range symmetric matrix
$\Gamma$ is positive definite and hence defines a Gaussian covariance.

\begin{proposition}[Window estimates from bounds on the non-flat matrix]
\label{ch06:windows}
Assume $cI\le\mathsf A\le CI$ and an exponentially weighted Schur bound
$\sup_i\sum_{i'}e^{\tau d(i,i')}|\mathsf A_{ii'}|\le C$, with the analogous
column bound. Assume bounded coordinate multiplicity per coarse site.
Then, for some $\eta,c'>0$, uniformly in the active set, period, and
principal compression,
\begin{equation}
 \|P-\mathsf A^{-1}\|_\eta+
 \|\Gamma-\mathsf A^{-1}\|_\eta\le C'e^{-c'B}.
 \label{ch06:window-error}
\end{equation}
Here $\|\cdot\|_\eta$ is the maximum of the exponentially weighted row and
column sums. For sufficiently large $B$, $c_1I\le\Gamma\le C_1I$.
Fixed-order derivatives obey corresponding bounds if the matrix $\mathsf A$
has them, with all windows and retained coordinate sets held fixed.
\end{proposition}
\begin{proof}
The weighted conjugation argument of Section~\ref{ch05:inverse} gives
exponential inverse bounds for $\mathsf A$ and all its principal
compressions. The constants are uniform because a principal compression
preserves the gap and the weighted bound. The column $Pe_i$ solves
$\mathsf APe_i=e_i$ inside $\Omega_i$. Outside that window its residual
crosses a distance at least $B/60$ from $i$. The convolution of two
exponentially decaying kernels therefore gives
$\|\mathsf AP-I\|_\eta\le C'e^{-c'B}$ after decreasing the decay rate.
Multiplication by $\mathsf A^{-1}$ proves the first estimate in
\eqref{ch06:window-error}; adjunction proves the second. Operator norm
is bounded by Schur norm, so comparison with $\mathsf A^{-1}\ge C^{-1}I$
proves positivity.

For derivatives, use $D\mathsf A^{-1}=-\mathsf A^{-1}(D\mathsf A)
\mathsf A^{-1}$ and its iterates, also on the compressions. Every term in
a differentiated residual still crosses the window boundary. Polynomial
spatial weights are moved along a kernel using
$1+d(o,i')\le(1+d(o,i))(1+d(i,i'))$ and absorbed by a smaller exponential
rate. These operations leave the factor $e^{-c'B}$ at any fixed order.
\end{proof}

\subsection{The inhomogeneous term and derivatives on a fixed sector}

The inhomogeneous term $b_Y$ in \eqref{ch06:actual-matrix} uses a bounded extension of the
frozen logarithms. This detail is essential. Let $v_Y$ be the patched,
smoothly clipped frozen displacement, extended by zero off the frozen
coordinates, with $|v_Y|_\infty\le2d$ and $d\le r t_f$, for fixed $r$.
If $b(V)$ is the value of the components at $\operatorname{rep}V$ and
$\mathsf D_{\rm all}(V)$ is its derivative on all non-tree coordinates,
including the frozen ones, the inhomogeneous term is
\begin{equation}
 b_Y=b(V)+\mathsf D_{\rm all}(V)v_Y,\qquad
 f_Y=\mathsf D^*b_Y .
 \label{ch06:clipped-source}
\end{equation}
The weights of the components are included in both terms. The original integrand is
preserved by putting the difference from this affine square into the
local correction factors. No estimate below replaces an arbitrary
frozen Haar logarithm by a small logarithm.

Here is why the inhomogeneous term is uniformly small. A retained component has an active
coordinate outside $A^{[10B]}$. Its whole
link support has radius at most $B/50$; hence its base point, endpoints, and fixed
macro-boxes have more than $9B$ clearance for large $B$. All coarse
variables $r_p^c$, $n_e$ used by the component are therefore small. In the common local gauge,
the four-link curvature formula, the absolute path-coefficient bounds,
and the smooth branch of the averaging operation give
\[
 |b(V)|_\infty\le C_b t_f,\qquad
 \|\mathsf D_{\rm all}(V)\|_{\infty\to\infty}\le C_a.
\]
The comparison term is extended without its large-field indicator when the field is
varied. Thus \eqref{ch06:clipped-source} implies
$|b_Y|_\infty\le(C_b+2rC_a)t_f$, uniformly over the frozen Haar variables.
The column bound for $\mathsf D$ now gives
\begin{equation}
 |f_Y|_\infty+|m_Y|_\infty\le C t_f,\qquad
 |\mathsf A m_Y+f_Y|_\infty\le C t_f e^{-cB}.
 \label{ch06:actual-forcing}
\end{equation}
The second inequality uses
$\mathsf A m_Y+f_Y=(I-\mathsf A P^*)f_Y$ and the window estimate.
In particular $g^{-1}(\mathsf A m_Y+f_Y)$ remains exponentially small
for $g=\beta^{-1/2}$ on the coupling band: $g^{-1}t_f=M_f^{\rm cut}$
up to fixed constants grows slower than $e^{\epsilon B}$.

For clarity, the corresponding derivative assertion has a different
scale. For any fixed allowed order $k$, use dimensionless variation
parameters whose coarse link increments and clipped frozen displacements
have derivatives bounded by $C_k t_f$. Keep the choice of retained components, the spanning trees, and
windows fixed. Differentiating a contour variable contributes a fixed power of its
length. The exponential path bound absorbs that power after decreasing
its decay rate. A derivative with base point $o$ can occur only on a contour variable
whose dependence set contains $o$; its distance to the component is bounded
by the length of the contour plus the fixed radius of the averaging operation. Consequently the
same argument bounds the derivatives of $\mathsf D$ in exponentially
weighted row and column norms, uniformly in volume. Products give the
bounds for $\mathsf A=\mathsf D^*\mathsf D$. All fixed positive-order
derivatives of $b_Y,f_Y,m_Y$ are $O(t_f)$ in these directions of size $O(t_f)$,
and the derivatives of the residual $\mathsf Am_Y+f_Y$ retain the factor $e^{-c_kB}$.
Unscaled derivatives instead carry the corresponding powers of $t_f^{-1}$.

These derivatives exist on a common neighborhood, not just on the
set satisfying the large-field indicator inequalities. Indeed the formulas for the components with the event indicator omitted
extend smoothly there. If their first derivative norm is at most $C$,
a ball of radius $c_L/[2(1+C)]$ preserves the gap $c_L/2$ of $\mathsf A$.
Every principal compression remains invertible on that ball. Inverse
differentiation and the proof of Proposition~\ref{ch06:windows} therefore
apply with the retained components fixed even if nearby fields change their large-field indicators.
The radius need not be uniform in arbitrary unnormalized coordinates.

\subsection{The exact integral and its normalization}

Let $n$ be the real dimension of the active coordinate space. Write $g=\beta^{-1/2}$, and let $\gamma_\Gamma$ denote the centered
Gaussian measure on $\RR^n$ with positive definite covariance $\Gamma$.
For any $m\in\RR^n$ and any $F$ integrable against the weight on the left,
the change of variables $u=m+gz$ gives
\begin{align}
 \int_{\RR^n}e^{-\beta Q_Y(u)}F(u)\,du
 &=\mathcal C_Y\int F(m+gz)
       \exp\bigl(\eta\cdot z+\tfrac12 z^*(\Gamma^{-1}-\mathsf A)z\bigr)
                                      \,d\gamma_\Gamma(z),\notag\\
 \mathcal C_Y&=(2\pi g^2)^{n/2}(\det\Gamma)^{1/2}e^{-\beta Q_Y(m)},\notag\\
 \eta&=-g^{-1}(\mathsf Am+f_Y).
 \label{ch06:gaussian-identity}
\end{align}
This is the usual finite-dimensional completion of squares with an
arbitrary reference covariance. We use it as an identity, not an
approximation. In particular $m=m_Y$ and \eqref{ch06:windows-def} do not
remove the remaining linear or quadratic term in the exponent.

In a sector integral $F$ includes the chart indicators, the Haar density
of every active link, the probability weights of the auxiliary index, incoming activities, and all
remaining action terms. It is extended by zero off its original chart.
The frozen variables are subsequently integrated with probability Haar
measure $dY$. When $\Gamma$ depends on $Y$, the resulting reference law is
\[
 dY\,d\gamma_{\Gamma(Y)}(z).
\]
It is a conditional Gaussian law, generally not a product of an independent
Haar field and a fixed Gaussian. The positive scalar $\mathcal C_Y$ must
be retained with the Haar Jacobians and relative determinants. A change of
tree coordinates preserves this combined integral, not its determinant
factor in isolation. Zero-dimensional active sectors have determinant one
and no Gaussian variables.

\section{Differentiation on a common dependence set}

The cutoffs and large-field indicators in the sector definition are fixed Borel conditions.
Their indicators are not differentiated. Even on a fixed sector an
incoming large-field activity can be merely measurable. Differentiating
it as a function of its fine links would be unjustified. A local change
of variables instead keeps those physical fine links fixed.

Write an active physical link as $U_e=\exp(u_e)(\operatorname{rep}V)_e$.
At a base field $V_0$, set $u=m_Y(V_0)+gz$. Within the logarithmic chart of the active links, where these logarithms are defined (its radius $d_*$ is fixed below), define
\begin{align}
 u_V(e)&=\log\!\left(\exp(u_e)(\operatorname{rep}V_0)_e
                                  (\operatorname{rep}V)_e^{-1}\right),\notag\\
 z_V&=g^{-1}(u_V-m_Y(V)).
 \label{ch06:protected-change}
\end{align}
Then $\exp(u_V(e))(\operatorname{rep}V)_e=U_e$ exactly. The following
extension makes this a legitimate change on the whole Gaussian space.

Choose a fixed logarithm radius $d_*>0$ and a smooth radial function
$\chi$ equal to one on $|u|\le d_*$ and zero on $|u|\ge2d_*$. For a
group element $M$ sufficiently near the identity put
\[
 F_M(u)=u+\psi(M,u),\qquad
 \psi(M,u)=\chi(u)\bigl(\log(\exp(u)M)-u\bigr).
\]
The logarithm is used only on the support of $\chi$. At $M=1$,
$\psi=0$ and $D_u\psi=0$. Compactness of that support permits one common
neighborhood of $1$ on which
$|\psi|<d_*/2$ and $\|D_u\psi\|<1/2$.
For every target $v$, the map $u\mapsto v-\psi(M,u)$ is a contraction
of the complete Euclidean space, so $F_M$ is bijective.
Its derivative is invertible, and the inverse function theorem gives a
smooth inverse. Its determinant is positive by continuity from $M=1$.
All fixed parameter derivatives of $\psi$ and of this determinant are
bounded uniformly in $u$, since they vanish outside a fixed compact set.
These are the usual finite-dimensional contraction and inverse-function
arguments.

Apply this map separately on each active bond with
$M_e(V)=(\operatorname{rep}V_0)_e(\operatorname{rep}V)_e^{-1}$, and set
\[
 T_Vz=g^{-1}\bigl(F_{M(V)}(m_Y(V_0)+gz)-m_Y(V)\bigr).
\]
The extension of $m_Y$ described above, without the event indicator, is used off the valid
sector. The map is globally smooth and invertible, and its Jacobian is
the product of the positive bond Jacobians; the factors $g$ cancel.
Call a factor \emph{protected} when it is cut off inside this logarithmic chart. If a protected factor is supported where $|u_V(e)|\le8a_*t_f$
and $8a_*t_f\le d_*/2$, then $|u_e|<d_*$ by the displacement bound.
There $F_M$ equals the logarithmic map and holds the physical link fixed.
Nested cutoffs, with the smaller support inside $|u_V|\le4a_*t_f$,
leave a strict margin. Thus a protected measurable factor is constant
on a neighborhood of the base parameter whenever it contributes; it
does not have to be differentiated. Cutoff remainders (defined in the next section) must likewise
vanish on a neighborhood, rather than only at the base point.

Use one coordinate change on the union $I$ of all coordinates read by the
product under consideration. If $\Gamma_I$ is its marginal covariance,
the density ratio contributed by that change is
\begin{equation}
 \mathcal R_I(V,z)=|\det D_z z_V|\,
 \frac{(\det\Gamma_I(V))^{-1/2}
        e^{-z_V^*\Gamma_I(V)^{-1}z_V/2}}
      {(\det\Gamma_I(V_0))^{-1/2}
        e^{-z^*\Gamma_I(V_0)^{-1}z/2}}.
 \label{ch06:density-ratio}
\end{equation}
The marginal is the principal covariance block $\Gamma_I$, whose inverse
is generally not a principal block of $\Gamma^{-1}$. Applying separate
changes to overlapping factors would not reproduce this common ratio.

The familiar derivative formulas
\[
 D\log\det\Gamma_I=\tr(\Gamma_I^{-1}D\Gamma_I),\qquad
 D\Gamma_I^{-1}=-\Gamma_I^{-1}(D\Gamma_I)\Gamma_I^{-1}
\]
show that the logarithmic derivative of the Gaussian part is
\[
 -\tfrac12\tr(\Gamma_I^{-1}\dot\Gamma_I)
 +\tfrac12z^*\Gamma_I^{-1}\dot\Gamma_I\Gamma_I^{-1}z
 -z^*\Gamma_I^{-1}\dot z.
\]
Higher fixed derivatives are polynomials in the coordinates, with
coefficients controlled by the window estimates and the coordinate
change. Ordinary Gaussian moment formulas, in particular
Isserlis' formula~\cite{Isserlis1918}, then give polynomial bounds in $|I|$
for fixed derivative order. The trace contributes a factor proportional to
$|I|$; it is not uniformly bounded independently of the dependence set.
For local factors with $s$ base points one typically has $|I|\le CB^4s$.

A factor's dependence set includes every cutoff condition, every contour of the averaging operation,
window matrix entry and frozen coordinate needed to evaluate it. This set
is different from the exclusion support $Y_i$ in
Definition~\ref{ch04:activities}. Disjoint exclusion supports imply no
independence. Even disjoint coordinate sets give conditional Gaussian
independence only when their cross covariance vanishes. Integration over
shared frozen variables can destroy that independence after conditioning
is removed.

\subsection{Differentiation under the common Gaussian law}

There is a useful distinction between the bounds used to justify
differentiation and those used to estimate the resulting derivatives.
The first may depend on the fixed finite volume. Since $\psi$ has
compact support and the frozen logarithms are clipped,
$T_Vz-z$ and its fixed parameter derivatives are bounded uniformly in
$z$ and in the frozen variables on one neighborhood. If
$\Gamma_I(V_0)^{-1}\ge\lambda I$, shrink that neighborhood so
$\Gamma_I(V)^{-1}\ge\lambda I/2$. The Gaussian density ratio and its
first $k$ derivatives are then bounded by
\[
 C\,p_{\Gamma_I(V_0)}(z)^{-1}
       e^{-\lambda|z|^2/8}(1+|z|)^{2k},
\]
with a finite constant $C$. Multiplication by the reference Gaussian
density makes this integrable. The Jacobian has bounded derivatives.
For the local correction factors, the endpoint cutoffs ensure bounded
arguments whenever the factor is nonzero. Hence their derivatives are
also bounded on that neighborhood. Dominated differentiation follows.
For a quadratic correction this last assertion requires both endpoints
of every nonzero matrix term to carry the protecting cutoff; it is not
a consequence of a cutoff at only one endpoint.

At the base point $T_{V_0}z=z$ and the ratio is one. Its differentiated
logarithm is a polynomial of degree at most two in $z$ and the
derivatives of $T_Vz$, and repeated exponential differentiation gives
a fixed-degree polynomial bound. Matrix products are estimated in row
and column norms; a trace contributes a factor $|I|$. Therefore, if $s$ counts the
base points of the factors and $|I|\le C B^4s$, there are fixed integers $r_k,m_k$
such that
\begin{equation}
 |D^j\mathcal R_I(V_0,z)|
 \le A_k(\widehat\beta_f)(1+s)^{m_k}
             (1+|I|+\sum_{i\in I}|z_i|)^{r_k},\quad j\le k,
 \label{ch06:base-scores}
\end{equation}
where $A_k$ is subpower, meaning
$A_k(x)\le C_\epsilon x^\epsilon$ eventually for each $\epsilon>0$.
In this statement the normalized variation derivatives of $g^{-1}m_Y$
and of $g^{-1}\psi$ must have subpower bounds; these follow from
\eqref{ch06:actual-forcing} and $t_f/g\asymp M_f^{\rm cut}$.
The relevant marginal inverses have uniform operator bounds; when
row bounds are used, a polynomial factor in $|I|$ suffices and is
already allowed in \eqref{ch06:base-scores}. No assertion about
independence of the coordinates is involved.

\section{Simultaneous estimates for large fluctuation-field events}

A cutoff remainder, a cutoff factor minus one, is bounded but need not be pointwise small. Its
smallness comes from a large fluctuation-field event. An event label specifies one
particular cutoff and its base point; repeated appearances of that event label
must be counted once. Suppose each remainder with event label $w$ vanishes
unless
\begin{equation}
 \mathcal E_w=\{\exists i:\ d(i,w)\le 2B,
                                   \ |z_i|\ge T\},\qquad
 T=\xi t_f/(2g),
 \label{ch06:tail-event}
\end{equation}
for fixed $\xi>0$. Verifying this implication uses the mean $m_Y$, the
derivative of the averaging operation, and the cutoff factors. A covariance estimate by itself cannot
prove it.

\begin{lemma}[Joint tails with bounded local multiplicity]
\label{ch06:joint-tails}
Let $z$ be zero-mean Gaussian with covariance $C\le M I$. Suppose an
event label has at most $N$ possible witness coordinates, and each coordinate
can witness at most $Q$ distinct event labels. For $r$ distinct event labels,
if $T\ge M/2+HQ$ with $H>0$, then
\[
 \mathbb P\left(\bigcap_{w=1}^r\mathcal E_w\right)
                                      \le(2N e^{-H})^r.
\]
No independence between the events is assumed.
\end{lemma}
\begin{proof}
This is the exponential Markov bound for a Gaussian, with the local
multiplicities kept explicit. For any distinct coordinate set $J$ and
signs $\sigma_i$, the variance of $\sum_{i\in J}\sigma_i z_i$ is at most
$M|J|$. Therefore
\[
 \mathbb P\{|z_i|\ge T\ (i\in J)\}
            \le2^{|J|}\exp\bigl((M/2-T)|J|\bigr).
\]
Assign one witness to each event label. There are at most $N^r$ assignments.
For the distinct coordinates in an assignment, $|J|\ge r/Q$ and
$|J|\le r$. The displayed estimate is at most $2^re^{-Hr}$.
Summing assignments proves the assertion.
\end{proof}

For the four-dimensional lattice, bounded multiplicity of coordinates
and event labels per coarse site gives $N,Q\le C_jL^d(1+B)^4$ for a fixed
finite power $d$. For $\beta$ in the coupling band,
$t_f/g\asymp M_f^{\rm cut}=\exp((\log\log\widehat\beta_f)^2)$.
It grows faster than every power of $\log\widehat\beta_f$. Thus, for any
fixed $\tau>0$, one can arrange
$H\ge2\tau\log\widehat\beta_f+\log(2N)$ and obtain
\begin{equation}
 \mathbb P\left(\bigcap_{w=1}^r\mathcal E_w\right)
                                      \le\widehat\beta_f^{-2\tau r}.
 \label{ch06:power-tail}
\end{equation}
The same statement holds uniformly under conditioning on the frozen variables if the
covariance bound and local vanishing assertion are uniform in that
conditioning.

\subsection{The deterministic estimate behind a cutoff remainder}

To distinguish the analytic issue from the probability estimate, write
$B_p(U,V)$ for a fine curvature component or averaging component minus its comparison term, and
$R_p(v)$ for its parallel-transported affine approximation on the same local
patch. On the complement of the large fluctuation-field event, $v=m_Y+gz$ has size $O(t_f)$.
Taylor expansion on the common coarse gauge, with its polynomial-in-$B$
gauge bound, gives
\[
 |B_p(U,V)-R_p(v)|
 \le C B^a t_f^2+C B^a t_f e^{-\sigma R_*}=o(t_f).
\]
The first term contains the product-logarithm and parallel-transport errors;
the second is the omitted path tail. Here the fixed path-length truncation
$R_*$ is a positive fixed multiple of $B$, so its exponential dominates
all powers of the coupling. The row bound also gives
$|R_p(v)-R_p(m_Y)|\le C\xi t_f$ if all components $gz_i$ of the fluctuation field in its
dependence set obey $|gz_i|\le\xi t_f$.

It remains to bound the affine component at the mean, not just the mean:
\begin{equation}
 |R_p(m_Y)|\le C B^a t_f^2+C B^a t_f e^{-\sigma R_*},
 \label{ch06:mean-target-needed}
\end{equation}
or the corresponding bound after subtracting the prescribed nonzero
linear target. If that target is within the strict plateau margin, choose
$\xi$ so $C\xi$ is less than one third of the margin and increase the
lower bound on the reference coupling to absorb the two errors. The local factors of the two partition-of-unity weights (the weights given to the sets $A_f$ and $A_n$ by the global partition of unity of the first section, with local factors $\vartheta_f(r_p^f)$ or $1-\vartheta_f(r_p^f)$ and $\vartheta_n(n_e)$ or $1-\vartheta_n(n_e)$)
and the normalized weights of the auxiliary index are then equal to one. The same strict
inequalities persist on a parameter neighborhood and supply the
zero-germ assertion used in \eqref{ch06:pointwise-jets}.

There is a concrete comparison identity for the mean behind this estimate. Let $v$
be one candidate field on a common clean ball of radius $8B$, zero on
the tree, with $|v|_\infty\le C B^a t_f$. Suppose every affine component in
that ball satisfies $|\mathsf Dv+b_Y|\le\epsilon$, where
$\epsilon= C B^a t_f^2+C B^a t_f e^{-\sigma R_*}$.
Put $E=\mathsf AP-I$. Then, with inactive identity coordinates included,
\[
 m_Y-v=E^*v-P^*(\mathsf Av+f_Y),\qquad
 \mathsf Av+f_Y=\mathsf D^*(\mathsf Dv+b_Y)+\Pi_{\rm inactive}v.
\]
These are exact identities. For a target at distance greater than $30B$ from the location of $A$,
all the windows in this identity lie in the $8B$ ball about the target,
which contains no plaquette of $A_f$, edge of $A_n$ or site of $A_c$. The inactive term vanishes
there, since all non-tree coordinates are active and $v$ is zero on
the tree. Hence
\[
 |m_Y(i)-v(i)|\le C\epsilon+C e^{-cB}|v|_\infty
 \quad\text{within distance }B\text{ of the target}.
\]
The row bound transfers the same estimate to the affine component at the target.

\begin{proposition}[Cancellation for the affine mean]
\label{ch06:mean-cancellation}
Use the quadratic forms of the two consecutive scales as in
\eqref{ch06:flat-identities}, the representative $W_F$, in the block axial gauge, of the linearization of their critical configuration, and the comparison term
defined by parallel-transporting and truncating the untruncated flat maps $M_f,M_n,N$. Then
\eqref{ch06:mean-target-needed} holds uniformly at targets farther than
$30B$ from the location of $A$. The constants may depend on the fixed step.
\end{proposition}
\begin{proof}
Let $a=W_F h$ be the fine field that the linearization of the critical configuration, in its representative $W_F$ of \eqref{ch05:stationary-tangent}, assigns to a coarse
logarithmic field $h$, and set $w=X_c^{\rm lin}(h)$ and
$p=(X_f^{\rm lin}(a),Z^{\rm lin}(a,h))$. Write
$T w=(M_fw,M_nw)$. Quadratic Gaussian matching gives
$|p|_2^2=|w|_2^2$, with the unnormalized $\ell^2$ norms. The two identities
\eqref{ch06:flat-identities} give $T^*p=w$ and $T^*T+N^*N=I$.
Consequently
\[
 \|p-Tw\|_2^2
 =\|w\|_2^2-2\|w\|_2^2+\|Tw\|_2^2=-\|Nw\|_2^2.
\]
Both sides have opposite signs, so $p=Tw$ and $Nw=0$.
In particular every linear component minus its comparison term vanishes. This argument
also applies to each signed orientation with the image of the spanning trees under the symmetry;
it never assumes that the operator $W_F$ for rooted trees commutes with the symmetry.

Choose one coarse gauge on the clean $8B$ ball and let $h$ be its
logarithmic field, extended by zero outside. The local axial gauge
construction bounds $\|h\|_\infty\le C B^a t_f$.
The operator $W_F$ has a bounded row norm at fixed $L$:
the exponentially decaying kernel of the linearization is summed absolutely,
and the correction bringing it to the block axial gauge is a sum along finitely many tree
edges per block. Thus the relative candidate
$v=W_Fh-\operatorname{rep}h$ is zero on the tree and has
$\|v\|_\infty\le C B^a t_f$.

Compare a nonlinear component in the inner ball with its untruncated linear component.
Taylor's theorem for the four-link curvature and the fixed
branch of the averaging operation bounds the nonlinear remainder of the component by $C\|h\|_\infty^2$.
For a path term, differentiating its products adds powers of the
length, summable against the exponential path bound.
Parallel transports change by at most their length times
$C\|h\|_\infty$; multiplying by an input of the component contributes a further
$O(\|h\|_\infty)$. The same estimates apply to the comparison term, since
the comparison term uses exactly the parallel-transported untruncated maps appearing in $T$.
Its clipping is inactive on this clean link support.
Omitted paths contribute at most
$C B^a t_f e^{-\sigma R_*}$ after decreasing $\sigma$.

Finally the derivative in the fluctuation field is evaluated at $\operatorname{rep}V$,
whereas the parallel-transported flat derivative is evaluated at zero before
parallel transport. Their difference has row norm at most
$C B^a t_f$, by the same second-derivative and parallel-transport estimates.
Applied to $v$, it contributes $C B^{2a}t_f^2$. No equality of these two
matrices is used. The frozen part of the inhomogeneous term is absent on the clean inner ball:
all its non-tree coordinates are active. The exact linear cancellation
therefore gives
$|\mathsf Dv+b_Y|\le C B^{a'}t_f^2+
 C B^{a'}t_f e^{-\sigma R_*}$ simultaneously on every required component,
using the one field $h$. The comparison identity for the mean above proves
\eqref{ch06:mean-target-needed}.
\end{proof}

This conclusion uses the path realization, with its factorization, of the untruncated flat maps in
\eqref{ch06:flat-identities}. A different comparison term with the same numerical
norm bounds would not satisfy it. For event indicators far from the location of $A$ it establishes the
implication as follows. The plaquette part of the fine comparison term is at most
$C L^{-2}t_f$, and the comparison term of an averaging component at most $C t_f$.
Choose $L$ and the plateau width $W$ for the deviations $n_e$ so these bounds lie
strictly inside their respective plateau thresholds, as in the energy
estimate. The preceding proposition, the nonlinear Taylor error, and
the bound on the fluctuation field then leave the fine plaquette distance from the identity below $t_f/4$
and the deviation $n_e$ below $Wt_f$. The corresponding factors of the partition-of-unity weights and the
normalized weights of the auxiliary index there are equal to one. The bond increment is bounded
by $Ct_f$ on this event, so a sufficiently broad fixed chart cutoff is
also one. By convention, a cutoff factor other than a chart cutoff whose location (the projected corners of its plaquette, the endpoints of its edge, or its site) meets $A^{[300B]}$, and a chart cutoff on a link whose projected endpoints meet $A^{[304B]}$, is not written as one plus a cutoff remainder: it is kept whole as a factor of the sector, and its cutoff remainder is set to zero. Section~\ref{ch06:primitive-bounds} below imposes the corresponding conditions through distance $300B$ and retains the chart cutoffs through $304B$.\mgubi{Convention taken from Lean (YM.\allowbreak ConnectedActivityE6.\allowbreak Coat.\allowbreak mandatory, mayor, aura; definitions read by two agents) and existence/\allowbreak gaussian-warnings.\allowbreak tex; v3 did not state it. Not checked: that Lean's factor types are exactly v3's cutoff factors.} Every other cutoff remainder lies outside these neighborhoods and has these small-field neighborhoods.
Replacing $\xi$ by $\xi/2$ gives strict margins on a neighborhood and
hence the zero-germ support condition. This proves the implication
\eqref{ch06:tail-event} for these event indicators, with the fixed comparison term
and plateau choices just stated.

\subsection{Derivatives of products of local factors}
\label{ch06:bounded-jet-route}

The local expansion offers a direct way to combine differentiation
and rare events. It is useful even when no aggregate quadratic majorant
has been established. Let $F(V,z)$ be a product of local factors after the
single change of variables \eqref{ch06:protected-change} above. Suppose, at the base point, its derivatives
through order $k$ satisfy
\begin{equation}
 |D^jF(V_0,z)|\le A_k(x)(1+s)^{m_k}
       \mathbf1_{\mathcal E}(z)\prod_i a_i,\qquad j\le k,
 \quad x=\widehat\beta_f .
 \label{ch06:pointwise-jets}
\end{equation}
Here $\mathcal E$ is the intersection of the events for $r$ distinct
cutoff remainders. This support assertion concerns all derivatives:
each remainder must vanish identically on a parameter neighborhood
outside its event. Assume the local endpoint protections needed for
dominated differentiation in the preceding subsection.

Then the common Gaussian integral, including its density ratio, obeys
\begin{equation}
 \left|D^j\mathbb E_{\Gamma_I(V_0)}
                  [\mathcal R_I(V,z)F(V,z)]_{V=V_0}\right|
 \le \widetilde A_k(x)(1+s)^{\widetilde m_k}
                   x^{-\tau r}\prod_i a_i,\qquad j\le k,
 \label{ch06:actual-jet-integral}
\end{equation}
provided \eqref{ch06:power-tail} holds with exponent $2\tau r$.
The constants are uniform in the frozen variables, which can subsequently
be integrated against their Haar probability measure.

Here is the moment argument in full. If every Gaussian coordinate has
variance at most $M$, the elementary exponential estimate gives
$\mathbb E(1+|z_i|)^q\le2e^{Mq^2/2}$ for each positive integer $q$.
Writing $N=|I|$ and $F_I=1+N+\sum_I|z_i|$, one has
$F_I\le(1+2N)(1+\max_I|z_i|)$ and hence
\[
 \mathbb E F_I^q\le(1+2N)^q(1+2N e^{Mq^2/2}).
\]
This uses no independence. Apply Leibniz' rule,
\eqref{ch06:base-scores}, and \eqref{ch06:pointwise-jets}.
Cauchy--Schwarz gives
$\mathbb E[\mathbf1_{\mathcal E}F_I^q]
 \le\mathbb P(\mathcal E)^{1/2}(\mathbb E F_I^{2q})^{1/2}$.
Since $N\le CB^4s$, all remaining factors are polynomial in $s$ and
subpower in $x$, proving \eqref{ch06:actual-jet-integral}.

For its application, \eqref{ch06:pointwise-jets} must be verified for the
factors of the expansion, with their full dependence sets. Ordinary
exponentiated correction terms satisfy it once their logarithmic
derivatives have small amplitude: the partition formula for derivatives
of $e^h-1$ gives a fixed multiple of that amplitude when $|h|\le1$.
Large-field components use their smooth local extension only when the
entire required coarse neighborhood is small-field. Their physical
indicator is not differentiated. The residual Gaussian density, the exponential factor in \eqref{ch06:gaussian-identity}, is first
decomposed into the protected linear and quadratic correction
factors; it is not copied once for every activity. These are specific
conditions on the exact expansion, not automatic consequences of the
Gaussian gap.

\subsection{One majorant for the product}

Here $z$ has the conditional law $d\gamma_{\Gamma(Y)}$, with frozen links $Y$ integrated in probability Haar measure. The normalization $\mathcal C_Y$ and the residual exponential are both those of \eqref{ch06:gaussian-identity}.

The residual Gaussian density in \eqref{ch06:gaussian-identity}, the
Jacobian ratios, and differentiated factors produce unbounded functions
of $z$. Bounds on their separate integrals cannot simply be multiplied.
The correct estimate is under a single common Gaussian law.

For a decomposition that retains an unbounded exponential instead, the following alternative criterion is available. Suppose a finite product of
local factors has, after up to $k$ allowed derivatives, a majorant of the
form
\begin{equation}
 C_k(\widehat\beta_f)(1+s)^d\,P_k(z)
       \exp\bigl(\eta\cdot z+\tfrac12z^*Kz\bigr)
       \prod_i a_i\prod_w M_w\mathbf1_{\mathcal E_w},
 \label{ch06:majorant-form}
\end{equation}
where $s$ counts base points, $P_k$ is a fixed-degree polynomial in absolute
coordinates with controlled moments, $K\ge0$, and the product over $w$
uses distinct event labels. The numbers $a_i$ bound ordinary local factors;
$M_w$ bound the cutoff factors or the probability weights of the branches. Their number must
include all allowed local values of the auxiliary index. Assume
\begin{equation}
 \|\Gamma^{1/2}K\Gamma^{1/2}\|\le1/8,
 \qquad \tr(\Gamma^{1/2}K\Gamma^{1/2})
              +|\Gamma^{1/2}\eta|^2\le C\sum_i\lambda_i.
 \label{ch06:quadratic-control}
\end{equation}
These are bounds on the sum of all local corrections. Smallness of each
summand separately does not suffice when the number of factors grows.
Exponential kernel bounds and bounded overlap of their full dependence sets
are one way to obtain them. For the linear term, a useful estimate
uses weights $w_{ia}>0$ with $\sup_a\sum_iw_{ia}\le C$; then
$|\sum_i\eta_i(a)|^2\le C\sum_i|\eta_i(a)|^2/w_{ia}$.
It controls the cross terms in the common square.

\begin{proposition}[Common local integral bound]
\label{ch06:common-majorant}
Assume \eqref{ch06:majorant-form}--\eqref{ch06:quadratic-control}, with
uniform polynomial moment bounds for $P_k$ under the corresponding
tilted Gaussians. Assume the local vanishing implication
\eqref{ch06:tail-event} for the same coordinates and event indicators. Then the
absolute differentiated integral is at most
\[
 C'_k(\widehat\beta_f)(1+s)^{d'}
       e^{C'\sum_i\lambda_i}\prod_i a_i
                  \prod_{w\text{ distinct}}(M_w\widehat\beta_f^{-\tau}).
\]
The statement holds after integration of the frozen Haar variables when
all inputs are uniform in those variables. It requires no independence
between the local factors or the large fluctuation-field events.
\end{proposition}
\begin{proof}
The standard Gaussian quadratic-exponential formula, with
$\widetilde K=\Gamma^{1/2}K\Gamma^{1/2}$, is
\[
 \mathbb E_\Gamma e^{\eta\cdot z+z^*Kz/2}
 =\det(I-\widetilde K)^{-1/2}
   \exp\left(\tfrac12\eta^*\Gamma^{1/2}
           (I-\widetilde K)^{-1}\Gamma^{1/2}\eta\right).
\]
Use it also with doubled exponent to control the square of the majorant.
The strict norm inequality in \eqref{ch06:quadratic-control} bounds the tilted
covariance and mean; Gaussian polynomial moments remain controlled.
Cauchy--Schwarz and \eqref{ch06:power-tail} give the asserted bound.
\end{proof}

Finally, Proposition~\ref{ch06:energy} contributes
$e^{-\kappa\beta t_f^2|A|}$. Throughout the coupling band,
$\beta t_f^2\asymp(M_f^{\rm cut})^2$. Any verified
logarithmic losses per element of $A$ of size
$O(B^d(1+\log\widehat\beta_f)^d)+\beta O(B^dt_f^3)$ are
$o(\beta t_f^2)$. They can be included while retaining
$e^{-\kappa'\beta t_f^2|A|}$ for any fixed $0<\kappa'<\kappa$ once the
lower bound on the reference coupling is increased. This last comparison concerns losses
actually assigned per element of $A$. It cannot absorb an unexplained
volume-dependent normalization.

\section{Weighted bounds for the elementary local factors}
\label{ch06:primitive-bounds}

We now retain the individual correction terms, including repeated paths,
instead of estimating only their sum as an operator. Fix a derivative
order $k$, a support exponent $\alpha>0$, and a multiplier $h$ such that the $hB$-neighborhoods are
large enough to contain every window and protected variation used above.
For a support $S$, let $c_B(S)$ be the least number of radius-$B$
balls whose union is connected and contains $S$. For each factor
$\lambda$, let $S_\lambda$ contain all sites on which it depends.
We bound its nonnegative amplitude $a_\lambda$ by requiring
\[
 \sup_x\sum_{\lambda:x\in S_\lambda}
               e^{\alpha c_B(S_\lambda)}a_\lambda\le\epsilon.
\]
The statement concerns the positive derivative orders available in the
incoming path bounds. Additional fixed polynomial factors in $B$ or
$M_f^{\rm cut}$ are subpower in $x=\widehat\beta_f$.

\subsection{Gaussian correction series}

Put $D=\mathsf A\Gamma-I$. The window estimates give exponentially small
weighted entry and fixed derivative envelopes for $D$. For the
interpolating precision $\mathsf A_s=I+s(\mathsf A-I)$, $0\le s\le1$,
the same argument applies uniformly: its gap is at least
$\min(1,c_L)$ and its weighted bounds are uniform.
Let $P_s$ be its localized approximate inverse and $E_s=\mathsf A_sP_s-I$.
The exact inverse and trace expansions are
\[
 \mathsf A_s^{-1}=P_s\sum_{n\ge0}(-E_s)^n,\qquad
 \Gamma^{-1}-\mathsf A=\sum_{n\ge0}(-D)^{n+1}\mathsf A .
\]
All series converge in weighted matrix norm. Integrating
$\tr((\mathsf A-I)\mathsf A_s^{-1})$ gives $\log\det\mathsf A$.
The $n=0$ term is its local diagonal normalization; the remaining terms of the series
are $(\mathsf A-I)P_s(-E_s)^{n+1}$.
The second determinant series comes from
$\log\det(I+D)=\log\det\mathsf A+\log\det\Gamma$.
This determinant identity is legitimate even if $D$ is not symmetric:
$I+D=\mathsf A\Gamma$ is similar to a positive matrix, and the
convergent real logarithm series agrees with its positive determinant.

Each term is indexed by its ordered sequence of bonds, with the intermediate Lie algebra indices contracted in each matrix product. For each fixed derivative order $k$,
the corresponding derivative of the amplitude is bounded by
\[
 C_k(B)e^{-cB}\bigl(C_k(B)e^{-cB}\bigr)^n,
\]
with possibly a smaller positive $c$. The endpoint cutoffs in a quadratic term
bound each Gaussian endpoint and its normalized derivatives by
a subpower function. A remaining linear term has one such endpoint.
Inactive endpoints give a zero germ because both precision and
covariance have identity inactive blocks. Thus all these estimates hold
for the protected factors on the entire sample space.

For completeness, the sum over the supports of the terms is also small.
Consecutive bonds have distance at most $3B+2$.
There are at most $K=C_L(1+B)^4$ possible next bonds.
A term of length $n+u$, for fixed $u$, has support cover at most
$(n+u+1)s_h$, where $s_h$ depends only on the multiplier $h$.
If that support contains a prescribed site, its starting bond has
at most $C_L(1+B)^4(n+u+1)^4$ possible locations.
Fixed derivatives contribute an additional polynomial factor in $n+u$.
Absorb these polynomials by $C_k2^{(k+4)n}$. The containing-site sum
is consequently bounded by
\[
 C_k(B)e^{-cB}\sum_{n\ge0}
 \bigl(C_k(B)e^{-cB}2^{k+4}
                e^{\alpha s_h}K\bigr)^{n}.
\]
Its ratio is at most $1/2$ eventually. The sum is smaller than
$x^{-A}$ for every prescribed $A>0$. This proves smallness of the
individual determinant and inverse-correction terms with their full
supports, including repeated bonds, rather than merely a small
operator norm.

\subsection{Short local and inherited terms}

The remaining estimates for the elementary local factors can be displayed together. Each bound
below is a containing-site bound with full fixed $hB$-neighborhoods and fixed-order
normalized derivatives, on the coarse small-field neighborhoods needed
for those derivatives:
\[
\begin{array}{l|l}
 \text{family}&\text{bound}\\ \hline
 \text{nonlinear action remainders far from }A
       & C_k B^{a_k}\beta t_f^3\\
 \text{extracted mean squares far from }A
       & C_k B^{a_k}\beta t_f^4\\
 \text{local determinant differences}
       & C_k B^{a_k}t_f\\
 \text{local Haar logarithms}
       & C_k B^{a_k}t_f^2\\
 \text{linear Gaussian residual and correction terms}
       & x^{-A}\quad\text{for every fixed }A\\
 \text{inherited localized terms}
       & C_k B^{a_k}\delta_f\\
 \text{inherited $\beta$-independent quadratic terms}
       & C_k B^{a_k}t_f^2\Lambda_f\\
 \text{inherited cubic and quartic terms}
       & C_k B^{a_k}\beta(t_f^3+t_f^4)\\
 \text{omitted path tails}
       & C_k B^{a_k}\beta t_f^2 e^{-cR_*}.
\end{array}
\]
Here $\delta_f$ is the incoming bound on the localized terms and $\Lambda_f$
the bound on the coefficients of the $\beta$-independent quadratic term. For example
$\delta_f\le x^{-1/2}$ and $\Lambda_f\le x^{1/8}$ make every displayed
family small by a fixed negative power of $x$.
The factors from differentiation and counting can be absorbed by
reducing the decay exponent if necessary.

The first two lines use the affine square with the exact
derivative at the background field. Taylor's remainder is quadratic in the displacement;
multiplying by the base affine component, of size $O(t_f)$ on the protected
chart, gives $O(t_f^3)$. The square of the remainder is smaller.
The extracted mean square is $O(B^a t_f^4)$ by
Proposition~\ref{ch06:mean-cancellation}. For derivatives the directions
are the directions of size $O(t_f)$ used above; applying the
same finite Taylor formula with parameter derivatives gives the same
powers, with the fixed factors from the path bounds. No derivative of a large-field indicator is taken.

For the third line, compare non-flat and flat matrices in the common
small gauge on the local window. Their difference is $O(B^a t_f)$.
The resolvent identity for each principal compression transfers this
to the diagonal integral defining the local logarithmic normalization.
The window size and trace introduce only powers of $B$.
The fourth line follows from $\log J(u)=O(|u|^2)$ for the Haar density
in exponential coordinates and from the protected chart bound.
Every site count in an $hB$-neighborhood is polynomial in $B$.

For inherited terms there is a support issue in addition to these
pointwise powers. If an incoming support $X$ has fine connected covering number
$d=c_{B_f}(X)$, with $B_f\le2B$, then its projected $hB$-neighborhood satisfies
\[
 c_B(S)\le P_h+\frac{10P_h}{L}d,\qquad
 x_0\in S\Longrightarrow
 d(x_0,\pi o_X)\le(h+20)B(1+d),
\]
for a fixed $P_h$. To prove the first bound, traverse a spanning tree
of the fine covering centers, sample that walk every $LB$ steps, and
cover each projected $hB$-neighborhood by at most $P_h$ coarse balls.
The walk has length at most $10Bd$, giving the displayed coefficient.
The same walk proves the bound on the distance to the base point. Thus an incoming
base-point weight $e^{\sigma d}$ absorbs the coarse support weight and all
polynomial multiplicities provided
\[
 \alpha\,10P_h/L+4\epsilon<\sigma
\]
for some fixed $\epsilon>0$. Counting possible base points contributes
$C_L B^4(1+d)^4$; derivative factors contribute another fixed power of
$1+d$, absorbed by $e^{\epsilon d}$. Summation of the remaining
geometric series proves the inherited site bounds.
For path terms the analogous cost is their true length, including
winding and repetitions; the unused exponential path factor absorbs its
support and its omitted tail. This explains why a positive incoming
additional support decay, not merely an unweighted coefficient sum, is required.

\subsection{Components and assembly under one law}

All constituent factors keep their original full dependence sets, including cutoff and path conditions. Their exclusion supports, defined in Definition~\ref{ch04:activities}, determine the disjoint covering families, not independence; the parameter $\sigma$ in the weighted support sums is the exponential locality weight of \eqref{ch04:weighted-sums}.

The relevant energy is the remainder after Gaussian extraction,
not the local energy $\mathcal E_A$ alone. We specify the assignment to prevent a
positive affine square from being lost in this estimate.
Let $O$ be a component of $A$ with $n$ elements. Its assigned action
terms have base points within $200B$ of $O$; their \emph{core} consists of the
terms within $2B$. The remaining terms form an \emph{annulus} $Q$.
Form a graph on the subsets of $A$ by joining two sets whenever their
closed $500B$ neighborhoods intersect. Group the sets by the connected
components of this graph. Distinct components of this graph then have
disjoint $500B$ neighborhoods (separation greater than $1000B$ suffices),
so their $200B$ action and $304B$ Haar assignments cannot overlap.
The enlarged finite-range adjacency changes only the fixed polynomial
counting constants. This geometric disjointness does not assert
probabilistic independence.

For this component estimate impose the weak conditions of the partition of unity
through distance $300B$: every fine plaquette not in $A_f$ has distance from the identity at
most $t_f/2$, every deviation $n_e$ with $e\notin A_n$ is at most $2Wt_f$,
the prescribed coarse large-field indicators equal those of the configuration, and no coarse
large-field site occurs away from the sites of $A_c$. The retained values of the auxiliary index for the local gauge transformations
have positive probability weight. These are conditions on the factor itself;
they are stronger than the initial $5B$ consistency definition.
For an annular base point at distance greater than $2B$ from the location of $A$,
its required $B$ neighborhood and the smaller comparison patches
are consequently free of elements of $A$. All such reads fit inside
the retained $300B$ region. Smooth chart cutoff factors are retained
throughout $304B$.
Local determinant terms are assigned through $200B$, whereas Haar
terms are assigned through $304B$; these two sets are not identified.

Let $H$ be the fine curvature and averaging components extracted from the
annulus by the link-support rule of Section~\ref{ch06:sector}. The link support of a retained component has radius
at most $B/50$. A retained component meets an active coordinate outside
$A^{[10B]}$, so its entire link support lies outside
$A^{[(10-2/50)B]}$ and hence outside the $2B$ core.
We first work in one signed orientation, before its probability weight
is applied. Denote the derivative and the inhomogeneous term of its components by
$\mathsf D_\sigma,b_{Y,\sigma}$, including the plaquette and penalty
coefficients but excluding the square root of the orientation weight.
The global matrix in \eqref{ch06:actual-matrix} stacks
$\sqrt{\omega_\sigma}\mathsf D_\sigma$.
Suppress $\sigma$ temporarily and set
\[
 X_H(\widetilde u)=\tfrac12
       \sum_{j\in H}|(\mathsf D_\sigma\widetilde u_{\rm active}+b_{Y,\sigma})_j|^2.
\]
The orientation probability is not included in this half-square; it
will appear once in the average below. The coefficients of the components and
the frozen part of the displacement field $\widetilde u$ are included once: equivalently apply the all-coordinate derivative to
the full displacement field and add the value of the components at $\operatorname{rep}V$. The \emph{displacement field} $\widetilde u$
is the active increment, the clipped frozen logarithm on frozen
non-tree bonds, and zero on tree bonds. Where the retained cutoff factors are nonzero
it reconstructs the original fine field on every retained link support.
Indeed the clean annular fine plaquettes and deviations $n_e$,
together with the common-tree comparison, put each frozen mismatch
within its logarithm cutoff's inner plateau. On active bonds the
cutoff support gives $|\widetilde u|\le2w t_f$; take the fixed width $w$
larger than the frozen comparison constant. Both the tree field and $\operatorname{rep}V$
are identity on tree bonds. Thus this same displacement field, uniformly bounded
by $2w t_f$, may be used in the affine approximation; an unrelated
logarithm or an unrestricted Gaussian sample cannot be substituted.

Write $W(Q)$ for the sum of the six local energy terms
\eqref{ch06:six-terms} on the annulus, with the same fine clip threshold
as the core. Then
\begin{equation}
 W(Q)-X_H(\widetilde u)\ge-\epsilon(x)n,\qquad
 \epsilon(x)\le C B^a t_f^3+C B^a t_f^2e^{-cR_*}
                         =o(t_f^2).
 \label{ch06:annulus-extraction}
\end{equation}
To verify the inequality, work in one clean local gauge about each
annular action term. Its nonlinear component minus its comparison term has norm
$M\le C B^a t_f$ and differs from its affine approximation
by $E\le C B^a t_f^2+C B^a t_f e^{-cR_*}$.
Taylor expansion gives this for the derivative at $\operatorname{rep}V$
directly; comparing first to parallel-transported flat components adds an error of
the same order. Hence the difference of the two half-squares is at
most $E(2M+E)/2$. The energies $E_f$ and $E_n$ equal the squares of their
nonlinear components minus comparison terms, plus the nonnegative difference between the squared distance from the identity and
the squared curvature representative. If a component is not retained, no affine square is subtracted,
so only its nonnegative square remains.
The coarse counterterm contributes $O(t_f^4)$.
The linear term of the complement component vanishes by the cancellation in
Proposition~\ref{ch06:mean-cancellation}; its nonlinear square is
$O(B^a t_f^4)$. The two signed mismatch terms are a comparison term of size
$O(t_f)$ paired with the quadratic Taylor/parallel-transport errors, hence
$O(B^a t_f^3)+O(B^a t_f^2e^{-cR_*})$.
There are at most $C B^4n$ such terms, proving
\eqref{ch06:annulus-extraction}.
The clean annular fields are below the fine threshold, so the fine
clip is its quadratic branch. The retained indicators of the relaxed regularity condition also equal
one: within their nontrivial range the patch is clean, and beyond
that range their permitted bound exceeds the uniform bound on the distance from the identity in the group. This checks the formula for the components used in the argument.

For precision, let $\sigma$ range over the signed lattice orientations
with their probability weights $\omega_\sigma$, transforming the field,
spanning trees, auxiliary index, index triple, and components together. Every term inside the
following average is the unweighted single-orientation quantity just
defined. Thus each orientation probability is counted exactly once. The energy assigned to this component
after Gaussian extraction is exactly
\begin{equation}
 \mathcal E_O^{\rm rem}
 =\operatorname{Av}_{\sigma}\bigl(
       \mathcal E_{\sigma O}+W(Q_\sigma)
       -X_{H_\sigma}(\widetilde u_\sigma)+C_\sigma(m_Y)\bigr),
 \qquad C_\sigma(m_Y)\ge0 .
 \label{ch06:component-extracted-energy}
\end{equation}
This is a finite identity between action terms: every near action term is assigned once;
a retained fine curvature or averaging component subtracts its affine square;
its extracted mean square is retained with positive sign.
Components that are not retained subtract zero. The mean-square sum is $C_\sigma$.
The near action terms partition into the core and annulus, so
\eqref{ch06:component-extracted-energy} and
\eqref{ch06:annulus-extraction}, followed by Proposition~\ref{ch06:energy},
give
\[
 \mathcal E_O^{\rm rem}\ge(\kappa t_f^2-\epsilon(x))n
                         \ge\tfrac12\kappa t_f^2 n
\]
at large coupling. This is the needed bound on the extracted energy.
Dropping $X_H$ before proving the annular estimate would be invalid.

The remaining local normalization, Haar, auxiliary-index, and frozen-coordinate
contributions are bounded by $C B^a(1+\log x)^a n=o(\beta t_f^2)n$.
The component amplitude consequently satisfies
\[
 a_{\rm comp}\le\exp\bigl(-\kappa'\beta t_f^2n\bigr)
\]
for a fixed positive $\kappa'$ smaller than the retained energy
coefficient. Fixed derivatives of its smooth extension on a fully
small-field coarse neighborhood add polynomial factors per component
and are absorbed by a further decrease of $\kappa'$. The physical Borel
indicator is held fixed. This assertion is where the retained cutoff factors are nonzero,
under the precise core/annulus assignment just specified; where a
cutoff factor kept inside the component vanishes, the full component factor is zero.
Connected index triples with prescribed root have at most
$\bigl(C(1+B)^a\bigr)^{n}$ possibilities: list them by a rooted spanning tree,
using their finite-range adjacency. This entropy and
$e^{\alpha c_B(S)}\le e^{C\alpha n}$ are absorbed by the same
remaining positive energy bound. A component's physical Borel indicator is held fixed;
the derivative assertion concerns its smooth extension on the stated
small-field coarse neighborhood.

Finally replace each small logarithm $f_\lambda$ by
$e^{f_\lambda}-1$. For $|f_\lambda|\le1$, the exponential partition
formula preserves its amplitude up to a constant depending on $k$.
The previous bounds furnish the pointwise product jets
\eqref{ch06:pointwise-jets}; derivative factors polynomial in the total
support size are absorbed by a small further decrease of the support exponent $\alpha$.
Cutoff remainders carry the joint-event factors proved above.
Equation~\eqref{ch06:actual-jet-integral} then gives a product majorant
under the one conditional Gaussian law, followed by the frozen
Haar probability integral. This is the local integration estimate
required before the connected expansion. Incoming large-field
activities retain their given amplitudes and participate in this same
product bound; they are not claimed to become pointwise small by
Gaussian integration.

\subsection{Comparison with the linear local integral}
\label{ch06:linear-comparison}

The estimates also give the small parameter needed to compare a local
nonlinear integral with its linear Gaussian counterpart. Fix the same
spanning trees, coordinate set, windows, and clean coarse field throughout the
comparison. Write $\Gamma_{\rm lin}$ for the window covariance formed
from the corresponding flat components, and $v$ for the
linear relative field $W_Fh-\operatorname{rep}h$ of the proof of Proposition~\ref{ch06:mean-cancellation} on the common clean neighborhood.
Taylor expansion of the nonlinear components and the principal resolvent identity
give, on those windows,
\[
 \|\Gamma-\Gamma_{\rm lin}\|\le C B^a t_f,\qquad
 |m_Y-v|_\infty\le C B^a t_f^2+C B^a t_f e^{-cB}.
\]
The second estimate is the comparison for the mean proved above; the first
is valid in weighted row norms, or with a harmless local polynomial
factor in an unweighted matrix norm. Both have fixed-order versions
under the same available path derivative bounds. In particular the
mean shift in Gaussian units satisfies
\[
 g^{-1}|m_Y-v|_\infty
       \le C B^a g(M_f^{\rm cut})^2+
                         C B^a M_f^{\rm cut}e^{-cB}.
\]
Thus it is of size $x^{-1/2+o(1)}$, rather than merely bounded.

The physical coordinate comparison uses the same block axial chart.
Its nonlinear map and its first-order Taylor map differ by
\[
 C B^a(t_f+g|z|)^2
       \le C'B^a(t_f^2+g^2|z|^2)
\]
on their protected common chart. Dividing by the fine small-field
radius gives $C B^a(t_f+(g^2/t_f)|z|^2)$.
All its fixed Gaussian moments are $x^{-1/2+o(1)}$ on a region
with polynomial-in-$B$ dimension. This is a Taylor estimate with
polynomial moment control, not a uniform bound on unbounded samples;
the chart-complement events are controlled separately.

For an interpolation of the positive covariances and of the centered
means, differentiation of its common Gaussian density ratio therefore
has local polynomial moments bounded by $x^{-1/2+o(1)}$, with the
appropriate polynomial factor in the size of the dependence set.
Indeed the logarithmic derivative is the trace, quadratic, and linear
score displayed after \eqref{ch06:density-ratio}; each coefficient
contains one of the covariance or normalized mean differences just
estimated. Integrating this derivative compares the two laws.
For an inherited activity with only $k$ available derivatives, this
comparison must be implemented by holding its physical fine links fixed
with the change of variables \eqref{ch06:protected-change}. Differentiate the common
density, Jacobian, smooth cutoff factors, and at most the available
$k$ activity derivatives. Taylor expansion of
$D^k E(U_{\rm nonlinear})-D^k E(U_{\rm linear})$ would require an
unavailable $(k+1)$st derivative and is not justified here.
The estimate above for the change of variables $T_V$ supplies bounds on
$T_V$ and its density scores; it is not permission to lose an extra
derivative of the incoming activity.
This comparison applies on a local region where all the windows lie in
the small-field region. It does not compare the Gaussian law of a sector whose index triple has nonempty $A_f$ or $A_n$ with the law of the sector with $A_f=A_n=\emptyset$,
does not replace $W_F$ by the representative in a different gauge,
and does not identify separately chosen event indicators. Such replacements need
their own tail or support decomposition.

\section{What must be checked for a concrete step}

The local integral estimate now follows from the gap of the non-flat matrix $\mathsf A$, the clipped
inhomogeneous term, cancellation of the comparison terms under the factorization hypothesis, strict margins in the event inequalities, and the
bounds for the elementary local factors above. Its hypotheses are the compatible untruncated
quadratic identities of the two consecutive scales, the rooted averaging operation, exponential path bounds
through the prescribed derivative order, and the incoming activity bounds
with a strictly larger support exponent than the one used in the conclusion. These are incoming data, not estimates
that can be created by choosing the schedule of Chapter~\ref{ch10:trajectory}. All matrices,
comparison terms, spanning trees, values of the auxiliary index, and frozen inhomogeneous terms in the argument belong to
that same pair of scales.

The chapter controls local amplitudes and their integrals. Reconstruction
of the exact block density additionally requires assigning every
nonlinear component, determinant contribution, Haar factor, and event indicator exactly
once. In particular the two near sets for local determinant terms and
Haar terms need not coincide, and the flat determinant scalar is
extracted only once for the whole family of spanning trees. Those identities, and the
subsequent connected regrouping, are distinct from the analytic bounds
proved here.
